\begin{abstract}
Training a nonlinear support vector machine (SVM) at scale is expensive
because decomposition solvers repeatedly generate kernel columns, update a
global gradient state, and execute a sequential working-set optimizer. This
report studies these costs on an NVIDIA RTX PRO 6000 Blackwell GPU using the
RBF C-SVC implementation in cuML 26.06 as the starting point. We redesign
kernel generation and reuse, then introduce a \emph{coherent-FP16} execution
path in which the persistent kernel cache and the local SMO subproblem consume
the same low-precision kernel representation. On EMBER-65K, coherent FP16
reduces median fit time from \SI{0.486954}{s} to \SI{0.360385}{s} relative to a
same-framework, equal-column-capacity direct FP32-cache control, a
\textbf{1.351$\times$} speedup over seven balanced fresh-process runs (95\%
bootstrap confidence interval: 1.263$\times$--1.403$\times$). The comparison
holds logical cache capacity fixed, not cache bytes. On all 200,000 official
EMBER test samples, a separately matched stock-model comparison changes
accuracy from 0.830980 to 0.830895 and ROC AUC from 0.93525055 to 0.93516107;
McNemar's test detects no statistically significant accuracy difference in
this experiment ($p=0.231645$). The coherent implementation does not satisfy
the original FP32-kernel KKT audit and must not be described as an exact FP32
drop-in solver.

After kernel acceleration, the serial SMO block solver becomes the dominant
GPU kernel, accounting for 40.3\% of the measured steady solver phase. We
evaluate whole-working-set projected and Newton methods, fixed-depth block-QP
updates, sampled candidate selection, geometry-aware pair packing,
stale-candidate prediction, exact-pair replay, triangular batching, and
speculative overlap. Real working-set captures and a 2,896-step exact GPU
trace show that these methods either introduce excessive dependent solves,
consume additional local kernel rows, increase global kernel-column work, or
retain a substantial serial step-and-update floor. The strongest
decision-path-aware block-solver micro-optimization provides only a
1.0376$\times$ paired geometric-mean end-to-end speedup. The resulting design
boundary is a work-depth constraint: a successful parallel SMO mechanism must
increase optimization progress per kernel data item while reducing dependent
stages and outer full-data scans.
\end{abstract}
